% Generated draft: requires mathematical and source review.
\documentclass[11pt]{article}
\usepackage[margin=25mm]{geometry}
\usepackage{fontspec}
\setmainfont{Latin Modern Roman}
\usepackage{amsmath,amssymb,mathtools,unicode-math}
\setmathfont{Latin Modern Math}
\usepackage{enumitem,needspace,xurl,hyperref}
\setlength{\parindent}{0pt}
\setlength{\parskip}{5pt}
\setlength{\emergencystretch}{3em}
\newcommand{\N}{\mathbb{N}}
\newcommand{\Z}{\mathbb{Z}}
\newcommand{\Q}{\mathbb{Q}}
\newcommand{\R}{\mathbb{R}}
\newcommand{\C}{\mathbb{C}}
\newcommand{\F}{\mathbb{F}}
\newcommand{\A}{\mathbb{A}}
\newcommand{\PP}{\mathbb{P}}
\newcommand{\E}{\mathbb{E}}
\newcommand{\eps}{\varepsilon}
\newcommand{\dd}{\,\mathrm{d}}
\begin{document}
\section*{UnsolvedMath ID: 30004878}
\textbf{Source problem code:} catalogue entry
\subsection{Definitions and mathematical statement}
The supplied record proposes a relation between poles of normalized R-matrices and first extension groups, but does not supply enough data to formulate a uniquely determined conjecture.

To describe the proposed target schematically, suppose that a coefficient field \(k\), a quantum loop algebra, a family \(\mathcal K\) of Kirillov--Reshetikhin modules with specified spectral parameters, an associated generalized preprojective algebra \(\Pi\), and an appropriate \(k\)-linear module category \(\mathcal C\) have been fixed. Suppose also that a specified construction assigns a generic-kernel module \(G(V)\in\mathcal C\) to each \(V\in\mathcal K\), and that a specified normalization and spectral-parameter convention give a meromorphic R-matrix \(R^{\mathrm{norm}}_{V,W}(z)\). These are required choices, not data determined by the supplied excerpt.

When the local coefficient functions of this R-matrix are defined and have at most a finite-order pole at \(z=1\), write
\[
 p_1\!\left(R^{\mathrm{norm}}_{V,W}\right)
 :=\min\left\{m\in\mathbb Z_{\geq0}:
 (z-1)^mR^{\mathrm{norm}}_{V,W}(z)
 \text{ is regular at }z=1\right\}.
\]
Thus this convention assigns pole order zero when the R-matrix is regular at \(1\). The excerpt does not specify this convention or the local realization needed to apply it. The group \(\operatorname{Ext}^1_{\mathcal C}(A,B)\) denotes extension classes represented by short exact sequences \(0\to B\to E\to A\to0\) in the specified category.

With these unresolved choices understood, the clean statement suggests the following candidate question:
\[
 \text{For every ordered pair }V,W\in\mathcal K,\qquad
 p_1\!\left(R^{\mathrm{norm}}_{V,W}\right)
 \stackrel{?}{=}
 \dim_k\operatorname{Ext}^1_{\mathcal C}\!\left(G(V),G(W)\right).
\]
This display is only a schematic rendering of the supplied clean statement. The field, algebraic type, admissible modules and parameter restrictions, normalization, generic-kernel construction, extension category, finiteness hypotheses, and intended orientation of the extension group remain unspecified. In particular, the record does not determine whether extensions are graded, ungraded, or aggregated over grading shifts.

The original extracted text begins by saying that an omitted conjecture can be checked whenever its unspecified left-hand side is known. It then discusses algebras denoted \(\Pi\) and a corrupted symbol, a \(\mathbb Z\)-grading, vertex simples \(S_i\), and modules \(E_i\) described as maximal indecomposable iterated self-extensions of \(S_i\) in the graded category. It asserts that a graded Euler--Poincar\'e pairing \(\langle E_i,S_j\rangle_q\) is a well-defined formal Laurent series expressible using a \(q\)-Cartan matrix \(C(q)\). It further describes graded submodules \(\bar I_i\) of vertex injectives satisfying \(\langle E_j,\bar I_i\rangle_q=\delta_{ij}\) and having graded dimensions in \(\mathbb Z_{\geq0}[q^{\pm1}]\). A subsequent formula labelled Theorem 1, attributed to reference [3, Theorem A], is corrupted. These are contextual assertions of the supplied fragment; they neither specify the omitted conjecture nor determine the displayed candidate equality.
\subsection{Short English statement}
The proposed question asks whether, for every admissible pair of Kirillov--Reshetikhin modules, the normalized R-matrix's pole order at \(z=1\) equals the dimension of a first extension group between associated generic-kernel modules. Essential definitions and hypotheses are missing, and the original fragment does not identify the conjecture. The displayed equality is therefore a provisional statement requiring source reconciliation and human mathematical review.
\subsection{Sources}
\begin{itemize}
\item Local UnsolvedMath registry: supplied source record for permanent ID \(30004878\), source problem code \texttt{catalogue entry}, title ``Pole Orders and Extensions of Kirillov--Reshetikhin Modules.'' The project identifies \texttt{data/unsolvedmath/problems.json} and \texttt{catalog/problems.json} as local registries; their contents were not independently inspected for this draft.
\item Supplied report citation: \emph{Mini-Workshop: Three Facets of R-Matrices (2022)}, Oberwolfach Reports, pp. 2791--2825, DOI \texttt{10.4171/owr/2021/51}. Supplied source URL and literature-source link, labelled ``OWR report'': \url{https://doi.org/10.4171/owr/2021/51}. The record separately gives proposed year 2021; no proposer is named.
\item Supplied literature assessment dated \texttt{2026-08-22}: the extracted sentence refers to an omitted conjecture and an unspecified left-hand side, preventing identification or verification of the mathematical claim; the record consequently labels it open. This is supplied dated curation, not an independently verified literature assessment.
\item The record also labels the statement \texttt{corrected\_verified} and asserts that it was checked against the official report and preserves its hypotheses, notation, and target. This supplied verification claim is unresolved against the limitations stated in the same record; no independent source visit or verification is claimed here.
\item The original fragment cites references [1] and [6] for appearances of an algebra in theoretical physics and quantum loop algebra representation theory, [6] for grading conventions, and [3, Theorem A] for Theorem 1. Bibliographic entries for these numbered references are absent from the supplied record.
\end{itemize}
\end{document}
